\documentclass[10pt,a4paper]{amsart}
\usepackage[margin=19mm]{geometry}
\usepackage{amsmath,amssymb,mathtools}
\usepackage[scr=rsfs,cal=euler]{mathalfa}
\usepackage{xfrac}
\usepackage[unicode,hidelinks,bookmarks=false]{hyperref}
\newcommand{\ie}{i.e.\ }
\newcommand{\cf}{cf.\ }
\newcommand{\resp}{respectively\ }
\newcommand{\Subsection}{\subsection}
\providecommand{\nobreakdash}{\nobreak\hbox{-}}
\setlength{\emergencystretch}{2em}
\multlinegap0pt
\usepackage{mathtools}
\usepackage{amssymb}
\usepackage{algorithm}\usepackage{algpseudocode}
\usepackage{xcolor}
\newtheorem{proposition}{Proposition}
\title{Randomized Subspace Nesterov Accelerated Gradient}
\author{Gaku Omiya, Pierre-Louis Poirion, Akiko Takeda}
\date{Source: arXiv:2605.00740v1 (CC BY 4.0)}
\begin{document}
\maketitle

\noindent\emph{Source and licence.} Randomized Subspace Nesterov Accelerated Gradient. Version arXiv:2605.00740v1 (CC BY 4.0), distributed by arXiv under the Creative Commons Attribution 4.0 International licence, http://creativecommons.org/licenses/by/4.0/, as recorded in the arXiv licence metadata for this exact version and shown on the abstract page https://arxiv.org/abs/2605.00740v1. Full licence notice: \texttt{licenses/arxiv-2605.00740-CC-BY-4.0.txt}.\par\smallskip

\noindent\textit{Editorial scope.} Counted unit: the article's proposition prop:strong-two-seq-obstruction-collapse (source0.tex lines 4387--4512). The article's complete original proof of it is delivered with it (source0.tex lines 4514--4789, one \texttt{proof} environment of 276 lines). No mathematics is written, repaired or re-derived: the statement and the proof are the article's own lines, in the article's own notation, and no retained line is substituted or re-flowed.  No further source-local context is needed: every cross-reference of every retained line resolves inside the two delivered spans.  Nothing was omitted from any retained span, so the package carries no omission record.  No retained line contains a citation, so the wrapper carries no bibliography and no bibliography entry.  No retained line contains a figure environment, an image-inclusion command or drawing code, so no image is delivered and the package declares no asset dependency.  Editorial formatting only, in addition: an AMS \texttt{amsart} preamble that restates the article's own declarations and macros used by the retained lines, namely the environments \texttt{proposition} on separate counters, together with the standard packages those lines need.  The retained environments print in this document as Proposition 1 in order of appearance, whereas the article prints the same items with the numbering of its own full document.  The complete native source of the article as distributed in the arXiv e-print for this exact version is bundled unchanged as \texttt{sources/199588/source0.tex}.
\begin{proposition}
\label{prop:strong-two-seq-obstruction-collapse}
Let \(f:\mathbb{R}^d\to\mathbb{R}\) be \(\mu\)-strongly convex and \(L\)-smooth, with \(\mu>0\).
Let \(\{P_k\}_{k\ge 0}\) be an i.i.d.\ sketch sequence satisfying
\[
\mathbb{E}[P_kP_k^\top]=I_d,
\qquad
P_k^\top P_k=\frac{d}{r}I_r
\]
for some \(1\le r\le d\).
Consider the direct sketched two-sequence recursion
\begin{equation}
x_{k+1}=y_k-\eta P_kP_k^\top \nabla f(y_k),
\qquad
y_{k+1}=x_{k+1}+\beta(x_{k+1}-x_k),
\label{eq:strong-direct-two-seq-collapse}
\end{equation}
where \(\eta>0\) and \(\beta\in[0,1)\) are constants. Define
\[
\theta\coloneqq \frac{1-\beta}{1+\beta}\in(0,1],
\qquad
\gamma\coloneqq \frac{\eta}{\theta},
\qquad
z_k\coloneqq \frac{(1+\theta)y_k-x_k}{\theta},
\qquad
u_k\coloneqq (1-\theta)z_k+\theta y_k,
\]
and
\[
\Phi_k\coloneqq f(x_k)-f^\star+\frac{\mu}{2}\|z_k-x^\star\|^2.
\]
Then
\begin{equation}
y_k=\frac{1}{1+\theta}x_k+\frac{\theta}{1+\theta}z_k,
\qquad
z_{k+1}=u_k-\gamma P_kP_k^\top \nabla f(y_k).
\label{eq:strong-hidden-two-seq-collapse}
\end{equation}

Writing \(g_k\coloneqq \nabla f(y_k)\), \(M_k\coloneqq P_kP_k^\top\), and
\(\mathcal F_k\coloneqq \sigma(P_0,\dots,P_{k-1})\), one has
\begin{align}
&\mathbb{E}\!\left[\Phi_{k+1}\middle|\mathcal F_k\right]-(1-\theta)\Phi_k
\notag\\
&\le
\theta\bigl(f(y_k)-f^\star\bigr)
+(1-\theta)\bigl(f(y_k)-f(x_k)\bigr)
-\mu\gamma\langle g_k,y_k-x^\star\rangle
+\mu\gamma\frac{1-\theta}{\theta}\langle g_k,x_k-y_k\rangle
\notag\\
&\qquad
+\frac{\mu\theta}{2}\|y_k-x^\star\|^2
-\frac{\mu(1-\theta)}{2\theta}\|x_k-y_k\|^2
+\left(
-\eta+\frac{d}{2r}(L\eta^2+\mu\gamma^2)
\right)\|g_k\|^2.
\label{eq:strong-precollapse-general}
\end{align}

If one wants the right-hand side of \eqref{eq:strong-precollapse-general} to collapse into the same two brackets as in the standard strongly-convex proof, namely
\begin{align}
&\mathbb{E}\!\left[\Phi_{k+1}\middle|\mathcal F_k\right]-(1-\theta)\Phi_k
\notag\\
&\le
\theta\Bigl(
f(y_k)-f^\star-\langle g_k,y_k-x^\star\rangle+\frac{\mu}{2}\|y_k-x^\star\|^2
\Bigr)
\notag\\
&\qquad
+(1-\theta)\Bigl(
f(y_k)-f(x_k)+\langle g_k,x_k-y_k\rangle
\Bigr)
\notag\\
&\qquad
-\frac{\mu(1-\theta)}{2\theta}\|x_k-y_k\|^2
\label{eq:strong-desired-collapse}
\end{align}
then matching the coefficients of the inner-product terms yields
\begin{equation}
\mu\gamma=\theta,
\qquad\text{that is,}\qquad
\gamma=\frac{\theta}{\mu},
\qquad\text{equivalently}\qquad
\eta=\frac{\theta^2}{\mu}.
\label{eq:strong-coupling-choice-collapse}
\end{equation}
Under \eqref{eq:strong-coupling-choice-collapse}, \eqref{eq:strong-precollapse-general} becomes
\begin{align}
&\mathbb{E}\!\left[\Phi_{k+1}\middle|\mathcal F_k\right]-(1-\theta)\Phi_k
\notag\\
&\le
\theta\Bigl(
f(y_k)-f^\star-\langle g_k,y_k-x^\star\rangle+\frac{\mu}{2}\|y_k-x^\star\|^2
\Bigr)
\notag\\
&\qquad
+(1-\theta)\Bigl(
f(y_k)-f(x_k)+\langle g_k,x_k-y_k\rangle
\Bigr)
\notag\\
&\qquad
-\frac{\mu(1-\theta)}{2\theta}\|x_k-y_k\|^2
+\frac{\theta^2}{\mu}
\left[
-1+\frac{d}{2r}\left(1+\frac{L\theta^2}{\mu}\right)
\right]\|g_k\|^2.
\label{eq:strong-collapse-final}
\end{align}

Therefore, in order for the remaining coefficient of \(\|g_k\|^2\) in
\eqref{eq:strong-collapse-final} to be nonpositive, it is necessary that
\begin{equation}
\frac{d}{r}\left(1+\frac{L\theta^2}{\mu}\right)\le 2.
\label{eq:strong-obstruction-condition-theta-collapse}
\end{equation}
Equivalently, since \(\eta=\theta^2/\mu\),
\begin{equation}
\frac{d}{r}(1+L\eta)\le 2.
\label{eq:strong-obstruction-condition-eta-collapse}
\end{equation}
In particular, if \(r\le d/2\), then \eqref{eq:strong-obstruction-condition-theta-collapse}
cannot hold for any \(\beta\in[0,1)\).
Therefore, outside a near-full-dimensional regime, the direct sketched analogue of the
classical strongly-convex two-sequence Nesterov template is incompatible with the
standard hidden-sequence proof mechanism.
\end{proposition}

\begin{proof}
Fix \(k\ge0\), and let
\[
\mathcal F_k\coloneqq \sigma(P_0,\dots,P_{k-1}),
\qquad
M_k\coloneqq P_kP_k^\top,
\qquad
g_k\coloneqq \nabla f(y_k).
\]

We first rewrite \eqref{eq:strong-direct-two-seq-collapse} in hidden-sequence form.
By the definition of \(z_k\),
\[
z_k=\frac{(1+\theta)y_k-x_k}{\theta},
\]
so that
\[
y_k=\frac{1}{1+\theta}x_k+\frac{\theta}{1+\theta}z_k.
\]
Also,
\[
z_{k+1}
=\frac{(1+\theta)y_{k+1}-x_{k+1}}{\theta}.
\]
Using
\[
y_{k+1}=x_{k+1}+\frac{1-\theta}{1+\theta}(x_{k+1}-x_k),
\]
we obtain
\begin{align*}
z_{k+1}
&=
\frac{(1+\theta)y_{k+1}-x_{k+1}}{\theta}\\
&=
\frac{1}{\theta}x_{k+1}-\frac{1-\theta}{\theta}x_k\\
&=
\frac{1}{\theta}(y_k-\eta M_k g_k)-\frac{1-\theta}{\theta}x_k.
\end{align*}
On the other hand,
\[
(1-\theta)z_k+\theta y_k
=
\frac{1-\theta}{\theta}\bigl((1+\theta)y_k-x_k\bigr)+\theta y_k
=
\frac{1}{\theta}y_k-\frac{1-\theta}{\theta}x_k.
\]
Hence
\[
z_{k+1}=(1-\theta)z_k+\theta y_k-\frac{\eta}{\theta}M_k g_k
=u_k-\gamma M_k g_k,
\]
which proves \eqref{eq:strong-hidden-two-seq-collapse}.

We next derive the one-step Lyapunov estimate for a general \(\gamma\), and only then identify
the value of \(\gamma\) for which the same collapse pattern as in the standard proof occurs.

From
\[
x_k-y_k=\theta(y_k-z_k),
\]
we obtain
\[
z_k-x^\star
=
y_k-x^\star-\frac{1}{\theta}(x_k-y_k).
\]
Also, since
\[
u_k=(1-\theta)z_k+\theta y_k
=
y_k-\frac{1-\theta}{\theta}(x_k-y_k),
\]
the update \(z_{k+1}=u_k-\gamma M_k g_k\) yields
\[
z_{k+1}-x^\star
=
y_k-x^\star-\frac{1-\theta}{\theta}(x_k-y_k)-\gamma M_k g_k.
\]
Therefore, a direct expansion gives
\begin{align}
&\|z_{k+1}-x^\star\|^2-(1-\theta)\|z_k-x^\star\|^2
\notag\\
&=
\theta\|y_k-x^\star\|^2
-\frac{1-\theta}{\theta}\|x_k-y_k\|^2
-2\gamma\langle M_k g_k,y_k-x^\star\rangle
+\frac{2\gamma(1-\theta)}{\theta}\langle M_k g_k,x_k-y_k\rangle
+\gamma^2 g_k^\top M_k^2 g_k.
\label{eq:strong-square-expansion-general}
\end{align}

Now
\[
M_k^2=P_k(P_k^\top P_k)P_k^\top=\frac{d}{r}M_k.
\]
Since \(P_k\) is independent of \(\mathcal F_k\) and \(\mathbb E[M_k]=I_d\), we have
\[
\mathbb E[M_k\mid\mathcal F_k]=I_d,
\qquad
\mathbb E[M_k^2\mid\mathcal F_k]=\frac{d}{r}I_d.
\]
Taking conditional expectation in \eqref{eq:strong-square-expansion-general}, we obtain
\begin{align}
&\mathbb E\!\left[
\frac{\mu}{2}\|z_{k+1}-x^\star\|^2
\middle|\mathcal F_k
\right]
-\frac{(1-\theta)\mu}{2}\|z_k-x^\star\|^2
\notag\\
&=
\frac{\mu\theta}{2}\|y_k-x^\star\|^2
-\mu\gamma\langle g_k,y_k-x^\star\rangle
+\mu\gamma\frac{1-\theta}{\theta}\langle g_k,x_k-y_k\rangle
-\frac{\mu(1-\theta)}{2\theta}\|x_k-y_k\|^2
+\frac{d}{2r}\mu\gamma^2\|g_k\|^2.
\label{eq:strong-z-diff-general}
\end{align}

On the function-value side, by \(L\)-smoothness and
\[
x_{k+1}=y_k-\eta M_k g_k,
\]
we have
\[
f(x_{k+1})
\le
f(y_k)-\eta g_k^\top M_k g_k+\frac{L\eta^2}{2}g_k^\top M_k^2 g_k.
\]
Taking conditional expectation gives
\begin{equation}
\mathbb E[f(x_{k+1})\mid \mathcal F_k]
\le
f(y_k)-\eta\|g_k\|^2+\frac{d}{2r}L\eta^2\|g_k\|^2.
\label{eq:strong-smooth-step-general}
\end{equation}

We now expand \(\mathbb E[\Phi_{k+1}\mid\mathcal F_k]-(1-\theta)\Phi_k\):
\begin{align}
&\mathbb E\!\left[\Phi_{k+1}\middle|\mathcal F_k\right]-(1-\theta)\Phi_k
\notag\\
&=
\mathbb E[f(x_{k+1})\mid\mathcal F_k]-f^\star-(1-\theta)(f(x_k)-f^\star)
\notag\\
&\qquad
+\mathbb E\!\left[
\frac{\mu}{2}\|z_{k+1}-x^\star\|^2
\middle|\mathcal F_k
\right]
-\frac{(1-\theta)\mu}{2}\|z_k-x^\star\|^2
\notag\\
&=
\bigl(\mathbb E[f(x_{k+1})\mid\mathcal F_k]-f(y_k)\bigr)
+\theta(f(y_k)-f^\star)
+(1-\theta)(f(y_k)-f(x_k))
\notag\\
&\qquad
+\mathbb E\!\left[
\frac{\mu}{2}\|z_{k+1}-x^\star\|^2
\middle|\mathcal F_k
\right]
-\frac{(1-\theta)\mu}{2}\|z_k-x^\star\|^2.
\label{eq:strong-phi-decomp-general}
\end{align}
Substituting \eqref{eq:strong-z-diff-general} and \eqref{eq:strong-smooth-step-general} into
\eqref{eq:strong-phi-decomp-general}, we obtain
\begin{align}
&\mathbb E\!\left[\Phi_{k+1}\middle|\mathcal F_k\right]-(1-\theta)\Phi_k
\notag\\
&\le
\theta(f(y_k)-f^\star)
+(1-\theta)(f(y_k)-f(x_k))
-\mu\gamma\langle g_k,y_k-x^\star\rangle
+\mu\gamma\frac{1-\theta}{\theta}\langle g_k,x_k-y_k\rangle
\notag\\
&\qquad
+\frac{\mu\theta}{2}\|y_k-x^\star\|^2
-\frac{\mu(1-\theta)}{2\theta}\|x_k-y_k\|^2
+\left(
-\eta+\frac{d}{2r}(L\eta^2+\mu\gamma^2)
\right)\|g_k\|^2,
\end{align}
which is \eqref{eq:strong-precollapse-general}.

We now derive the value of \(\gamma\) for which the same collapse pattern as in the
standard proof occurs. In \eqref{eq:strong-precollapse-general}, the inner-product terms are
\[
-\mu\gamma\langle g_k,y_k-x^\star\rangle
+\mu\gamma\frac{1-\theta}{\theta}\langle g_k,x_k-y_k\rangle.
\]
In order to rewrite the inner-product terms in \eqref{eq:strong-precollapse-general}
with the same coefficients as in \eqref{eq:strong-desired-collapse}, one must impose
\[
\mu\gamma=\theta,
\qquad
\mu\gamma\frac{1-\theta}{\theta}=1-\theta.
\]
These two identities are equivalent, and therefore
\[
\gamma=\frac{\theta}{\mu}.
\]
Since \(\gamma=\eta/\theta\), this is equivalent to
\[
\eta=\frac{\theta^2}{\mu},
\]
which proves \eqref{eq:strong-coupling-choice-collapse}.

Substituting \(\gamma=\theta/\mu\) and \(\eta=\theta^2/\mu\) into
\eqref{eq:strong-precollapse-general}, we obtain
\begin{align}
&\mathbb E\!\left[\Phi_{k+1}\middle|\mathcal F_k\right]-(1-\theta)\Phi_k
\notag\\
&\le
\theta\Bigl(
f(y_k)-f^\star-\langle g_k,y_k-x^\star\rangle+\frac{\mu}{2}\|y_k-x^\star\|^2
\Bigr)
\notag\\
&\qquad
+(1-\theta)\Bigl(
f(y_k)-f(x_k)+\langle g_k,x_k-y_k\rangle
\Bigr)
\notag\\
&\qquad
-\frac{\mu(1-\theta)}{2\theta}\|x_k-y_k\|^2
+\left(
-\eta+\frac{d}{2r}(L\eta^2+\mu\gamma^2)
\right)\|g_k\|^2.
\label{eq:strong-after-matching}
\end{align}
Since
\[
-\eta+\frac{d}{2r}(L\eta^2+\mu\gamma^2)
=
-\frac{\theta^2}{\mu}
+\frac{d}{2r}
\left(
L\frac{\theta^4}{\mu^2}
+
\frac{\theta^2}{\mu}
\right)
=
\frac{\theta^2}{\mu}
\left[
-1+\frac{d}{2r}\left(1+\frac{L\theta^2}{\mu}\right)
\right],
\]
this yields \eqref{eq:strong-collapse-final}.

Finally, by \(\mu\)-strong convexity,
\[
f(y_k)-f^\star-\langle g_k,y_k-x^\star\rangle+\frac{\mu}{2}\|y_k-x^\star\|^2\le 0,
\]
and by convexity,
\[
f(y_k)-f(x_k)+\langle g_k,x_k-y_k\rangle\le 0.
\]
Thus, in order for the remaining coefficient of \(\|g_k\|^2\) in
\eqref{eq:strong-collapse-final} to be nonpositive, it is necessary that
\[
\frac{d}{r}\left(1+\frac{L\theta^2}{\mu}\right)\le 2,
\]
which is \eqref{eq:strong-obstruction-condition-theta-collapse}. The equivalent form
\eqref{eq:strong-obstruction-condition-eta-collapse} follows from \(\eta=\theta^2/\mu\).

Finally, if \(r\le d/2\), then \(d/r\ge2\), and hence for every \(\theta\in(0,1]\),
\[
\frac{d}{r}\left(1+\frac{L\theta^2}{\mu}\right)
\ge
2\left(1+\frac{L\theta^2}{\mu}\right)
>
2.
\]
Therefore \eqref{eq:strong-obstruction-condition-theta-collapse} cannot hold for any
\(\theta\in(0,1]\), and hence for any \(\beta\in[0,1)\).
This completes the proof.

\end{proof}

\end{document}
